%==============================================================
\section{Kết luận và hướng phát triển}
\label{sec:conclusion}
%==============================================================

\subsection{Kết luận}

Chuyên đề đã đi qua một lộ trình liền mạch từ khái niệm cơ bản nhất tới một bài toán nghiên cứu cụ
thể:
\begin{itemize}[noitemsep]
  \item \textbf{Nền tảng:} qubit là vector đơn vị trong $\mathbb{C}^2$ viết bằng ký hiệu Dirac; phép
    đo tuân theo quy tắc Born; relative phase quan sát được còn global phase thì không; hệ $n$ qubit
    có $2^n$ amplitude và có thể entangled.
  \item \textbf{Cổng lượng tử:} mọi cổng là ma trận unitary; cổng một qubit là phép quay trên Bloch
    sphere; CX là cổng hai qubit then chốt và là thước đo chi phí chính.
  \item \textbf{Thuật toán:} Deutsch--Jozsa, Bernstein--Vazirani và Grover dùng chung cơ chế phase
    kickback và interference; lợi thế được đo bằng số lần gọi oracle; teleportation và superdense
    coding dùng entanglement làm tài nguyên.
  \item \textbf{Compilation:} transpiler là một chuỗi pass trên \code{DAGCircuit} gồm sáu stage; ba
    bài toán con chính (layout, routing, synthesis) đều NP-hard và được giải bằng heuristic.
    Thực nghiệm cho thấy compiler ảnh hưởng mạnh tới chất lượng: \code{optimization\_level} nâng
    xác suất thành công từ 0.89 lên 0.94; layout tốt biến 3 SWAP thành 0; topology thay đổi số SWAP
    từ 75 xuống 25; \code{gridsynth} cần ít cổng $T$ hơn Solovay--Kitaev hơn 7 lần.
  \item \textbf{Compilation như MDP:} qubit routing được phát biểu thành MDP với state $(F,\pi)$,
    action là chọn cạnh SWAP, và ba kiểu reward không cùng chiều; môi trường \code{RoutingEnv} qua
    \code{check\_env} và kiểm tra cả tính hợp lệ lẫn tính đúng ở mọi episode; thang mốc cho thấy
    policy tham lam kém \code{SabreSwap(decay)} khoảng 25\%.
\end{itemize}

Bên cạnh kiến thức, chuyên đề rút ra một bài học phương pháp: trong lĩnh vực này, \textbf{sai lầm
nguy hiểm nhất là sai lầm im lặng}. Năm bẫy ở Bảng~\ref{tab:traps} đều không sinh lỗi mà chỉ sinh
số liệu sai, và nhiễu do seed ($\sim$9\%) hay do layout ($\sim$16\%) đủ lớn để biến một ``cải thiện''
thành nhiễu thuần. Vì vậy hai sản phẩm bàn giao quan trọng nhất của chuyên đề không phải là một
thuật toán mới, mà là một môi trường \textbf{đúng semantics} và một bộ \textbf{baseline trung thực}.

\subsection{Hạn chế}

\begin{itemize}[noitemsep]
  \item Mọi thí nghiệm chạy trên simulator và thiết bị giả lập (\code{GenericBackendV2}); chưa chạy
    trên phần cứng thật, nên reward \code{fidelity} mới dùng sai số giả lập.
  \item Quy mô thí nghiệm nhỏ (5--9 qubit, vài chục tới vài trăm cổng); kết luận có thể thay đổi ở
    quy mô của các thiết bị thật (100+ qubit).
  \item Chưa trình bày các thuật toán dựa trên Quantum Fourier Transform (phase estimation, Shor) và
    các thuật toán variational (VQE, QAOA).
  \item Chưa huấn luyện agent RL nào -- đây là chủ đích, vì môi trường và baseline phải được kiểm chứng
    trước.
\end{itemize}

\subsection{Hướng phát triển cho Chuyên đề nghiên cứu 2}

\begin{enumerate}[noitemsep]
  \item Huấn luyện \textbf{DQN} trên \code{RoutingEnv} với reward \code{swap} trên \code{line(5)} --
    bài toán nhỏ nhất có ý nghĩa. Mục tiêu đầu tiên là \textbf{vượt \code{GreedyPolicy}}, sau đó tiến
    tới \code{SabreSwap(decay)}.
  \item Dùng action masking trong mạng (gán logit của action bị mask bằng $-\infty$), sau đó thử
    \textbf{PPO}.
  \item So sánh bằng \code{bench.so\_voi\_baseline} trên cùng tập mạch, nhiều seed, báo cáo
    trung bình $\pm$ độ lệch chuẩn và tính đến mức nhiễu $\sim$9\%.
  \item Mở rộng dần topology: \code{line(5)} $\to$ \code{ring(5)} $\to$ \code{grid(3,3)} $\to$
    heavy-hex; thử reward \code{fidelity} với dữ liệu sai số thật; thử kiến trúc graph neural network
    để tổng quát hoá giữa các kích thước thiết bị \cite{riu2025zx}.
  \item Song song, thử bài toán \textbf{pass ordering} (action space nhỏ) để kiểm tra toàn bộ pipeline
    thí nghiệm, và tìm hiểu thêm QFT, phase estimation, VQE và mô phỏng có nhiễu bằng
    \code{qiskit\_aer.noise}.
\end{enumerate}
